\documentclass[11pt]{article}
\usepackage[margin=1in]{geometry}
\usepackage{amsmath,amssymb,amsthm}
\usepackage[hidelinks]{hyperref}
\setlength{\emergencystretch}{1em}
\hypersetup{pdftitle={Three anti-invariant classes on a nonintegrable four-manifold},
 pdfauthor={Alec Kriebel},pdfsubject={A counterexample to the anti-invariant integrability conjecture}}
\newtheorem{theorem}{Theorem}
\newcommand{\R}{\mathbb R}
\newcommand{\C}{\mathbb C}
\newcommand{\Z}{\mathbb Z}
\newcommand{\vol}{\mathrm{vol}}
\newcommand{\dd}{\mathrm d}
\newcommand{\Rea}{\operatorname{Re}}
\newcommand{\Ima}{\operatorname{Im}}
\title{Three anti-invariant classes on a\\ nonintegrable four-manifold}
\author{Alec Kriebel\\
 \small Independent researcher\\
 \small\href{https://orcid.org/0009-0001-9320-500X}{ORCID: 0009-0001-9320-500X}}
\date{October 16, 2026 \quad Version 1.0}
\begin{document}
\maketitle

\begin{abstract}
We construct smooth nonintegrable almost complex structures on the product
of a genus-two curve and an elliptic curve with at least three independent
real anti-invariant de Rham classes. The structures converge smoothly to the
product complex structure and equal it outside one coordinate ball.
The construction gives a negative answer to the unrestricted integrability
conjecture of Draghici--Li--Zhang. The separate generic-vanishing statement is
a prior theorem of Tan--Wang--Zhang--Zhu.
\end{abstract}

\section{Statement and context}
For an almost complex structure $J$ on a closed smooth four-manifold $M$, set
\[
 H_J^-=\{[\alpha]\in H^2_{\mathrm{dR}}(M;\R):
          \dd\alpha=0,\ \alpha(JX,JY)=-\alpha(X,Y)\},\qquad
 h_J^-=\dim_\R H_J^- .
\]
These groups and their cohomological setting are due to Li--Zhang and
Draghici--Li--Zhang \cite{LZ,DLZ08,DLZ}.
An early exact existence question appears in \cite[Question 3.27]{DLZ08}.
Conjecture 2.5 of \cite{DLZ} asserts that $h_J^-\geq3$ forces integrability.
It is repeated in Tomassini's contribution to Report 33/2020
\cite[p.~1687]{OWR}; the report's bibliographic year is 2020 and its EMS
publication date is September 14, 2021.

\begin{theorem}\label{thm}
Let $\Sigma$ be the smooth compact completion of $y^2=z^6-1$, let
$E=\C/(\Z+i\Z)$, and let $J_0$ be the product complex structure on
$M=\Sigma\times E$. There exist $\varepsilon_0>0$ and a smooth family
$J_\varepsilon$, $|\varepsilon|<\varepsilon_0$, such that
\begin{enumerate}
\item $J_\varepsilon=J_0$ outside a fixed coordinate ball and
      $J_\varepsilon\longrightarrow J_0$ in $C^\infty$;
\item three fixed linearly independent classes in $H^2_{\mathrm{dR}}(M;\R)$
      have closed $J_\varepsilon$-anti-invariant representatives;
\item $J_\varepsilon$ is nonintegrable whenever $\varepsilon\ne0$.
\end{enumerate}
In particular, $h_{J_\varepsilon}^-\geq3$ for $\varepsilon\ne0$.
\end{theorem}

The other conjecture in the source concerns generic vanishing. Tan, Wang,
Zhang and Zhu proved that $h_J^-=0$ is open and dense in the $C^\infty$
topology on all smooth almost complex structures on a closed four-manifold
\cite[Theorem 1.1]{TWZZ}. The special family above does not contradict
that theorem. We compute a lower bound, not the exact value of $h_J^-$.

\section{Construction and proof}
The polynomial $z^6-1$ has six simple roots. Its compact double cover of
$\mathbb{CP}^1$ is connected and has genus two by Riemann--Hurwitz.
The differentials
\[
 \eta_0=\frac{\dd z}{y},\qquad \eta_1=\frac{z\,\dd z}{y}
\]
extend holomorphically and are linearly independent. At a finite branch
point, $y$ is a coordinate and $\dd z/y=\dd y/(3z^5)$ has no pole.
At either point over infinity put $s=1/z$, $q=y/z^3$; then
$q^2=1-s^6$ and
\[
 \eta_0=-\frac{s\,\dd s}{q},\qquad \eta_1=-\frac{\dd s}{q}.
\]
Let $\dd w$ denote the global holomorphic one-form on $E$, and define
\[
 \Omega_j=\eta_j\wedge\dd w,\qquad
 \alpha=\Rea\Omega_0,\quad \beta=\Ima\Omega_0,\quad
 \gamma=\Rea\Omega_1 .
\]
These are closed and $J_0$-anti-invariant.

Their real cohomology classes are independent. Indeed, for real $a,b,c$,
put $H=(a-ib)\Omega_0+c\Omega_1$ and $F=\Rea H$.
Then $F=a\alpha+b\beta+c\gamma$ and
$F^2=\tfrac12 H\wedge\overline H$ is nonnegative in the complex orientation,
and positive wherever $H\ne0$. If $F$ were exact, Stokes' theorem would give
$\int_M F^2=0$, so $H=0$. Independence of $\eta_0,\eta_1$ gives $a=b=c=0$.

\paragraph{The simultaneous local model.}
Near $(z,y,w)=(0,i,0)$ choose the holomorphic branch $y=y(z)$ and a local
lift of $w$. The coordinates
\[
 \zeta=z,\qquad \xi=\frac{w}{y(z)}
\]
have nonzero Jacobian and satisfy
$\Omega_0=\dd\zeta\wedge\dd\xi$ and $\Omega_1=\zeta\Omega_0$.
Writing $\zeta=u+iv$ and $\xi=x+it$, we obtain
\begin{equation}\label{normal}
 \alpha=\dd u\wedge\dd x-\dd v\wedge\dd t,\qquad
 \beta=\dd u\wedge\dd t+\dd v\wedge\dd x,\qquad
 \gamma=u\alpha-v\beta .
\end{equation}
Here $t$ is a coordinate, and $\varepsilon$ is the deformation parameter.
Choose a relatively compact coordinate ball $B$ about the origin and a
smooth real cutoff $\rho\in C_c^\infty(B)$ equal to one near the origin.
Set
\[
 R=\frac{\varepsilon}{2}\rho x^2.
\]
The one-forms $R\,\dd v$ and $Rv\,\dd v$ extend smoothly by zero to $M$.
Define global closed forms
\begin{equation}\label{perturb}
 \alpha_\varepsilon=\alpha,\qquad
 \beta_\varepsilon=\beta+\dd(R\,\dd v),\qquad
 \gamma_\varepsilon=\gamma-\dd(Rv\,\dd v).
\end{equation}
Their cohomology classes are the three original independent classes.
The identity $\dd(Rv\,\dd v)=v\,\dd(R\,\dd v)$ holds for the full cutoff,
not only where $\rho=1$. Consequently,
\begin{equation}\label{third}
 \gamma_\varepsilon=u\alpha-v\beta_\varepsilon\quad\hbox{on }B .
\end{equation}
This simultaneous retention of the third closed class is the key relation.

\paragraph{Realizing the anti-invariant plane.}
Put $\vol=\dd u\wedge\dd v\wedge\dd x\wedge\dd t$. Subscripts on $R$
denote partial derivatives. Direct expansion gives
\[
 \beta_\varepsilon
 =\dd u\wedge\dd t+(1-R_x)\dd v\wedge\dd x
     -R_t\dd v\wedge\dd t+R_u\dd u\wedge\dd v,
\]
and hence
\begin{equation}\label{gram}
 \alpha^2=2\vol,\qquad
 \alpha\wedge\beta_\varepsilon=R_t\vol,\qquad
 \beta_\varepsilon^2=2(1-R_x)\vol .
\end{equation}
For sufficiently small $|\varepsilon|$ we have uniformly
$|R_x|\leq1/4$, $|R_t|\leq1$; these bounds follow because $\rho$ is fixed
and compactly supported. Thus
\[
 s=\frac{R_t}{2},\qquad D=1-R_x-\frac{R_t^2}{4}\geq\frac12 .
\]
Define on $B$
\[
 B_\varepsilon=\frac{\beta_\varepsilon-s\alpha}{\sqrt D},\qquad
 \Psi_\varepsilon=\alpha+iB_\varepsilon .
\]
Equation \eqref{gram} implies
\begin{equation}\label{volume}
 \Psi_\varepsilon^2=0,\qquad
 \Psi_\varepsilon\wedge\overline{\Psi_\varepsilon}=4\vol .
\end{equation}
In a complex vector space of dimension four a nonzero two-form with square
zero has rank two and is decomposable. Its contraction kernel therefore
has complex dimension two. The second identity in \eqref{volume} says that
this kernel and its conjugate are complementary. Declaring
\[
 T^{0,1}_{J_\varepsilon}
  =\{Z\in TM\otimes\C:\iota_Z\Psi_\varepsilon=0\}
\]
therefore defines a smooth real almost complex structure on $B$.
It agrees with $J_0$ at $\varepsilon=0$. This is the usual positive
two-plane description of an almost complex structure in dimension four
\cite{Donaldson}; the preceding kernel argument supplies the needed
realization directly.

On a collar of $\partial B$, all jets of $R$ vanish, so that
$s=0$, $D=1$, $\Psi_\varepsilon=\Omega_0$ and $J_\varepsilon=J_0$.
Extend $J_\varepsilon$ by $J_0$ outside $B$. The extension is smooth.
The fixed compact support and the uniform positive denominator imply
$J_\varepsilon\to J_0$ in every finite smooth norm, hence in $C^\infty$.
Zeros of $\Omega_0$ elsewhere on $M$ cause no difficulty: the kernel
construction is used only on $B$, where $\Omega_0$ is nonzero.

On $B$, $\alpha$ and $\beta_\varepsilon$ are real combinations of the real
and imaginary parts of $\Psi_\varepsilon$, so they are $J_\varepsilon$-anti-invariant.
Equation \eqref{third} gives the same conclusion for $\gamma_\varepsilon$.
Outside $B$ they are the original $J_0$-anti-invariant forms.
Equation \eqref{perturb} now proves the claimed three independent classes.

\paragraph{Nonintegrability.}
Where $\rho=1$ put $k=1-\varepsilon x>0$ and $k'=\partial_x k=-\varepsilon$.
There
\[
 \Psi_\varepsilon
  =(\dd u+i\sqrt{k}\,\dd v)\wedge
       (\dd x+i k^{-1/2}\dd t)=\phi_1\wedge\phi_2 .
\]
Thus $(\phi_1,\phi_2)$ is a $(1,0)$ coframe, and
\begin{equation}\label{nonint}
 \dd\phi_1=\frac{i k'}{2\sqrt{k}}\,\dd x\wedge\dd v,\qquad
 \dd\phi_1\wedge\phi_1\wedge\phi_2=\frac{k'}{2k}\vol .
\end{equation}
The last expression is nonzero for $\varepsilon\ne0$. For an integrable
structure, the differential of a $(1,0)$ form has no $(0,2)$ component, so
its wedge with a $(2,0)$ form vanishes in complex dimension two.
Equation \eqref{nonint} proves nonintegrability and completes the proof
of Theorem \ref{thm}.

\section{Comparison and limits}
The two-plane realization and closed-pair framework are established
ingredients \cite{Donaldson,BK}; the third-class identity \eqref{third}
is the specific step used here. Three independent cohomology classes do
not mean three pointwise independent two-forms: on $B$ their representatives
span a real plane, with the variable coefficients in \eqref{third}.

Two distinctions explain why familiar integrability criteria do not obstruct
this construction. First, on the core of $B$,
\[
 B_\varepsilon=k^{-1/2}\dd u\wedge\dd t+\sqrt{k}\,\dd v\wedge\dd x,\qquad
 \dd B_\varepsilon=-\frac{k'}{2k^{3/2}}\dd x\wedge\dd u\wedge\dd t\ne0 .
\]
A closed real pair need not remain closed after variable normalization.
In particular, $\Psi_\varepsilon$ is not a closed complex symplectic form;
the closed normalized-pair criterion \cite[Theorem 6]{BK} is inapplicable.
Second, on this core $J_0\partial_u=\partial_v$ whereas
$J_\varepsilon\partial_u=k^{-1/2}\partial_v$. A positive metric compatible
with both would give $g_{uu}=g_{vv}=k^{-1}g_{vv}$, impossible where
$k\ne1$. Thus no common compatible metric with this product $J_0$ exists.
No compatibility with a fixed symplectic form, spatial analyticity or exact
value of $h_{J_\varepsilon}^-$ is asserted.

Earlier compact four-dimensional nonintegrable examples attain $h_J^-=2$
\cite{DLZ,HT}. Hind--Tomassini's final arXiv version
\cite[Proposition 4.2]{HT} distinguishes this compact case from its
noncompact and real-six-dimensional larger examples; localized deformation
alone is consequently not a new ingredient.
Chen--Zhang's accepted manuscripts, I, \S6.3 \cite{CZI}, and II, \S5.1
\cite{CZII}, give closely related nonintegrable products $T^2\times S$. Their computed
plurigenera are different from the real closed anti-invariant cohomology
used here, and II explicitly leaves the latter calculation for those models
as an interesting question. We do not claim that no further argument about
those older models could supply such classes.

A bounded literature audit through October 16, 2026 found no verified earlier
qualifying counterexample. It does not establish firstness, novelty, present
openness or exhaustive absence. The supplementary audit records actual
reading and access limits: some journal-final comparisons were unavailable;
Chen--Zhang accepted versions were read; only selected author-uploaded
passages of the 2012 DLZ survey were accessible; and Geiges's 1996 paper
remained a bibliographic lead. Withdrawn claims \cite{LU,TWold} are not
used as theorems. These historical limits do not enter the explicit proof.

\paragraph{Verification and disclosures.}
AI tools were used extensively in the construction, drafting, literature
search, symbolic computations and adversarial checking. This research note
is unrefereed; no independent human peer review or formal verification is
claimed. The accompanying source and verification package preserves the
submitted proof and checks, independent audit reports and reproducible finite
algebra. Its integrity and computation checks do not replace the global
argument or certify historical priority. The author is responsible for the
content. This note is released under
\href{https://creativecommons.org/licenses/by/4.0/}{CC BY 4.0}.

\begin{thebibliography}{99}\small
\bibitem{LZ}
T.-J. Li and W. Zhang, Comparing tamed and compatible symplectic cones and
cohomological properties of almost complex manifolds,
\emph{Comm. Anal. Geom.} 17 (2009), 651--683.
\href{https://arxiv.org/abs/0708.2520}{arXiv:0708.2520}.

\bibitem{DLZ08}
T. Draghici, T.-J. Li and W. Zhang, Symplectic forms and cohomology
decomposition of almost complex 4-manifolds,
\href{https://arxiv.org/abs/0812.3680v1}{arXiv:0812.3680v1} (2008),
Question 3.27. A shortened version appeared in
\emph{Int. Math. Res. Not.} 2010, 1--17,
\href{https://doi.org/10.1093/imrn/rnp113}{doi:10.1093/imrn/rnp113};
the question number refers to v1.

\bibitem{DLZ}
T. Draghici, T.-J. Li and W. Zhang, On the $J$-anti-invariant cohomology of
almost complex 4-manifolds, \emph{Quart. J. Math.} 64 (2013), 83--111,
\href{https://doi.org/10.1093/qmath/har034}{doi:10.1093/qmath/har034}.
Also \href{https://arxiv.org/abs/1104.2511v1}{arXiv:1104.2511v1}.

\bibitem{OWR}
A. Tomassini (joint work with R.~K. Hind), On the cohomology of almost complex
manifolds, in \emph{Mini-Workshop: Almost Complex Geometry},
Report 33/2020, \emph{Oberwolfach Rep.} 17 (2020), 1686--1689;
published September 14, 2021,
\href{https://doi.org/10.4171/OWR/2020/33}{doi:10.4171/OWR/2020/33}.

\bibitem{TWZZ}
Q. Tan, H. Wang, Y. Zhang and P. Zhu, On cohomology of almost complex
4-manifolds, \emph{J. Geom. Anal.} 25 (2015), 1431--1443,
\href{https://doi.org/10.1007/s12220-014-9477-2}{doi:10.1007/s12220-014-9477-2}.
Also \href{https://arxiv.org/abs/1112.0768v3}{arXiv:1112.0768v3}.

\bibitem{Donaldson}
S.~K. Donaldson, Two-forms on four-manifolds and elliptic equations, in
\emph{Inspired by S.~S. Chern}, Nankai Tracts Math. 11 (2006), 153--172.
\href{https://arxiv.org/abs/math/0607083v1}{arXiv:math/0607083v1}.

\bibitem{BK}
G. Bande and D. Kotschick, The geometry of recursion operators,
\emph{Comm. Math. Phys.} 280 (2008), 737--749,
\href{https://doi.org/10.1007/s00220-008-0477-6}{doi:10.1007/s00220-008-0477-6}.
Also \href{https://arxiv.org/abs/math/0703348v2}{arXiv:math/0703348v2}.

\bibitem{HT}
R. Hind and A. Tomassini, On the anti-invariant cohomology of almost complex
manifolds, \emph{J. Geom. Anal.} 31 (2021), 4906--4922,
\href{https://doi.org/10.1007/s12220-020-00461-8}{doi:10.1007/s12220-020-00461-8}.
Numbering and comparisons here use
\href{https://arxiv.org/abs/1908.03016v3}{arXiv:1908.03016v3} (2020).

\bibitem{CZI}
H. Chen and W. Zhang, Kodaira dimensions of almost complex manifolds I,
\emph{Amer. J. Math.} 145 (2023), 477--514,
\href{https://doi.org/10.1353/ajm.2023.0011}{doi:10.1353/ajm.2023.0011}.
Comparisons use the
\href{https://wrap.warwick.ac.uk/id/eprint/164863/1/WRAP-Kodaira-dimensions-almost-complex-manifolds-I-Zhang-2022.pdf}{accepted manuscript},
\S6.3; see also
\href{https://arxiv.org/abs/1808.00885v2}{arXiv:1808.00885v2}.

\bibitem{CZII}
H. Chen and W. Zhang, Kodaira dimensions of almost complex manifolds II,
\emph{Comm. Anal. Geom.} 32 (2024), 751--790,
\href{https://doi.org/10.4310/CAG.241015030334}{doi:10.4310/CAG.241015030334}.
Comparisons use the
\href{https://wrap.warwick.ac.uk/id/eprint/163848/1/WRAP-Kodaira-dimensions-almost-complex-manifolds-II-Zhang-2022.pdf}{accepted manuscript},
\S5.1; see also
\href{https://arxiv.org/abs/2004.12825v1}{arXiv:2004.12825v1}.

\bibitem{LU}
M. Lejmi and M. Upmeier, An integrability theorem for almost-K\"ahler
structures using $J$-anti-invariant two-forms on four-manifolds,
\href{https://arxiv.org/abs/1507.00282}{arXiv:1507.00282},
withdrawn version 2 (September 3, 2015); not used.

\bibitem{TWold}
Q. Tan, H. Wang, Y. Zhang and P. Zhu,
Symplectic cohomology and the stability of $J$-anti-invariant cohomology,
\href{https://arxiv.org/abs/1307.1513}{arXiv:1307.1513},
withdrawn version 2 (November 13, 2014); not used.
\end{thebibliography}
\end{document}
